\documentclass[12pt]{article}
\usepackage{../../resource/shortex}
\usepackage{mathtools,amssymb}
\usepackage{enumitem}
\usepackage{datetime}
\usepackage{fancyhdr}
\usepackage{ifthen}
\usepackage{pifont}
\usepackage{mathrsfs}
\usepackage[normalem]{ulem}
%\RequirePackage[usenames,dvipsnames]{color}


%% showing solutions 
\newboolean{showsols}
\setboolean{showsols}{false}
\newcommand{\showsolutions}{\setboolean{showsols}{true}}
\newlength{\solspace}
\setlength{\solspace}{16em}
\newcommand{\nosolspace}{\setlength{\solspace}{0em}}

\newcommand{\solution}[2]{\ifthenelse{\boolean{showsols}}{{\par\textcolor{red}{Rubric criterion: #1}\newline\textcolor{blue}{#2}}}{\vspace{\solspace}}}
\newcommand{\npforprint}{\ifthenelse{\boolean{showsols}}{}{\newpage}}
\newcommand{\vsforprint}[1]{\ifthenelse{\boolean{showsols}}{}{\vspace{#1}}}

%% grading
\newcommand{\grades}[1]{~\newline\noindent\framebox[\textwidth][c]{E\quad \quad \quad \quad S \quad \quad \quad \quad R \quad \quad \quad \quad N}}

\newcommand{\questiongrades}{\noindent\emph{Each item is marked $\checkmark$, $\checkmark\!-$, or \ding{55}:}\begin{center}\emph{$\checkmark$ = capable  $~~$ $\checkmark\!-$ = mostly capable  $~~$ \ding{55} = not yet}\end{center} }

\newcommand{\triAssessments}{$\checkmark\quad\checkmark\!-\quad$\ding{55}\quad}
\newcommand{\biAssessments}{$\checkmark\quad$\ding{55}\quad}

%% line for writing name and BUID
\newcommand{\nameBUID}{\noindent\textbf{Name:}\underline{\phantom{XXXXXXXXXXXXXXXXXXX}}\hfill\textbf{BUID:}\underline{\phantom{XXXXXXXXXXXXX}}\newline}


%% formatting 
%\allowdisplaybreaks[3]
\renewcommand{\headrulewidth}{.4mm} % header line width
\renewcommand{\arraystretch}{1.25}


%% semester label, read from the "Semester:" line of important_dates.md
%% (path is relative to a document in materials/ModuleX/<folder>/)
\ExplSyntaxOn
\tl_new:N \g_teaching_semester_tl
\ior_new:N \g_teaching_dates_ior
\file_if_exist:nTF { ../../instructor_inputs/important_dates.md }
  {
    \ior_open:Nn \g_teaching_dates_ior { ../../instructor_inputs/important_dates.md }
    \ior_str_map_inline:Nn \g_teaching_dates_ior
      {
        \regex_extract_once:nnNT { \A \s* Semester: \s* (.*?) \s* \Z } {#1} \l_tmpa_seq
          {
            \tl_gset:Nx \g_teaching_semester_tl { \seq_item:Nn \l_tmpa_seq { 2 } }
            \ior_map_break:
          }
      }
    \ior_close:N \g_teaching_dates_ior
  }
  {
    \msg_new:nnn { teaching } { no-dates }
      { ../../instructor_inputs/important_dates.md~not~found;~semester~header~left~blank. }
    \msg_warning:nn { teaching } { no-dates }
  }
\cs_new:Npn \semester { \tl_use:N \g_teaching_semester_tl }
\ExplSyntaxOff

%% page style
\pagestyle{fancy}
\fancyhf{}
\rhead{\semester}
\lhead{CAS MA 213: Basic Statistics and Probability}
%\rfoot{\thepage}

%--------------------
\begin{document}

\begin{center}
\textbf{\Large Project 2, Part 1 \\ Lab Activity}
\end{center}

\begin{center}
\begin{tabular}{ll}
\toprule
\textbf{Item} & \textbf{Date} \\
\midrule
Lab P2\_1 (today) & Fri, Oct 23 \\
\textbf{Deliverable 1 (outline) due} & \textbf{Thu, Oct 29, 9:00 PM} \\
Lab 6 in class & Fri, Oct 30 \\
Lab P2\_2 in class & Fri, Nov 6 \\
\textbf{Deliverable 2 (results + R code) due} & \textbf{Thu, Nov 12, 9:00 PM} \\
Lab P2\_3 in class & Fri, Nov 13 \\
\textbf{Deliverable 3 (report + R code) due} & \textbf{Thu, Nov 19, 9:00 PM} \\
\bottomrule
\end{tabular}
\end{center}

\subsection*{Goal for Today}
Get your group set up to work efficiently on Project 2. You are \textbf{not} expected to finish Deliverable 1 in lab today --- the goal is a solid start.

\subsection*{Before You Start}
Copy the \href{https://docs.google.com/document/d/1J5nZ17STPfx3xI3VadGgF2XQOvTIQxo3MpW34h-n1bU/edit?usp=sharing}{P2\_D1\_Outline} Google doc and share it amongst your group members --- that is the document you will fill in as you go.

If you don't get to step 4 (assigning roles and scheduling your next meeting) with 5 minutes left, stop where you are and move straight to step 4 --- every group should leave lab today with roles assigned and a next meeting on the calendar.

\subsection*{Suggested Order}

\begin{enumerate}
    \item \textbf{Introduce group members ($\sim$3 min)}
    \begin{itemize}
        \item Go around the group, say your name and BU ID, and fill in the group members at the top of the outline (Title and Group Information).
        \item Identify your TF/LA mentor, and get their email address.
    \end{itemize}

    \item \textbf{Find a data set and questions that work together ($\sim$20 min)}

    Choosing a data set and choosing your questions are the same decision: a data set is only a good choice if it supports a categorical \emph{and} a numerical question that each match a test whose conditions are met. So work in a loop, and \textbf{don't fill in the outline yet} --- keep rough notes until you have settled.
    \begin{enumerate}
        \item \textbf{Pick a candidate data set} from the suggested sources below (or your own).
        \item \textbf{Load it in R} and look at the variables (\texttt{str()}, \texttt{summary()}, \texttt{table()}). If it won't load or is too messy, go back to (a).
        \item \textbf{Draft one categorical and one numerical question.} Look through the tests at the end of this sheet, find the ones that fit your variables, and use their examples as models.
        \item \textbf{Check the conditions in R.} For example, use \texttt{table()} to count successes and failures (or expected counts) in each group, check how many rows are in each group, and make a quick histogram to look for strong skew.
        \item \textbf{Decide:} if both questions work, you're done with this step. If one doesn't (no test fits, a condition fails, a group is tiny), go back to (c) with a different question, or to (a) with a different data set.
    \end{enumerate}
    \emph{Tip:} it is normal to go around this loop two or three times. It is much easier to switch data sets today than after Deliverable 1.

    \item \textbf{Fill in the outline ($\sim$12 min)}
    \begin{itemize}
        \item Data: source, what one row represents, and number of rows (1b); variables of interest (1c); the data quality table (1d).
        \item Analyses: for each of Sections 2 (categorical) and 3 (numerical), the research question and variables (a), a planned exploratory table or plot (b), and the planned test, hypotheses in words, and conditions (c).
        \item Now that you know your data and questions, write the project title (top of the outline) and background and motivation (1a), and add your data source to References (5).
    \end{itemize}

    \item \textbf{Assign roles by deliverable and plan your next meeting ($\sim$5 min)}
    \begin{itemize}
        \item Fill in the Roles by Deliverable grid (Roles and Timeline, at the end of the outline) and agree on a specific date/time (or async check-in method) for your next meeting, before Deliverable 1 is due.
    \end{itemize}
\end{enumerate}

\subsection*{Tip: Suggested Task List by Deliverable}

Use this as a menu when filling in the Roles by Deliverable grid.

\begin{description}[nosep, leftmargin=0pt, style=sameline, font=\normalfont]
    \item[\textbf{D1 --- Outline (write sections):}] title/group, 1a motivation, 1b data source, 1c variables, 1d data quality, 2 categorical plan, 3 numerical plan, 5 references
    \item[\textbf{D2 --- R analysis:}] data import and cleaning, categorical EDA (summary table + visualization), categorical inference (conditions, test statistic, p-value), numerical EDA (summary statistics table + visualization), numerical inference (conditions, test statistic, p-value)
    \item[\textbf{D2 --- Results doc:}] revising D1 sections from feedback, data cleaning (1d), categorical results (2b--2d), numerical results (3b--3d), draft conclusion (4)
    \item[\textbf{D2 --- Other:}] R code cleanup and assembly, data file, updating roles
    \item[\textbf{D3 --- Report sections:}] 1 introduction, 2 categorical analysis, 3 numerical analysis, 4 conclusion, 5 references, roles appendix
    \item[\textbf{D3 --- Other:}] revising R code from D2 feedback, proofreading and assembling the report
\end{description}

\subsection*{Possible Data Sources to Consider (Not Limited to These)}
You can click the links below to explore possible data sources.
\begin{itemize}
    \item \href{https://www.openintro.org/data/}{OpenIntro Datasets}
    \item \href{https://cran.r-project.org/web/packages/fivethirtyeight/vignettes/fivethirtyeight.html}{fivethirtyeight R Package}
    \item \href{https://stat.ethz.ch/R-manual/R-devel/library/datasets/html/00Index.html}{datasets R package}
    \item \href{https://www.kaggle.com/datasets}{Kaggle datasets}
    \item \href{https://data.gov/}{US government's open data}
\end{itemize}

\subsection*{Hypothesis Test Options and Their Conditions}
% Hypothesis test options for Project 2, shared by Project2.tex and
% P2_1_activity.tex. Each document supplies its own section heading.

Pick the tests that match your variables and research questions: at least one categorical test and at least one numerical test.
Every test has an \textbf{independence} condition: the observations should come from a random sample (or a randomized experiment), and if sampling without replacement, the sample should be less than 10\% of the population.

\subsubsection*{Categorical Data Analysis (choose at least one)}
\begin{itemize}
    \item \textbf{Single proportion} (Chapter 6.1): one categorical variable with two outcomes. Compare the proportion in one outcome to a hypothesized value $p_0$.\\
    \emph{Example:} Is the proportion of Titanic passengers who survived different from 50\%?\\
    \emph{Conditions:} independence; success-failure using the null value: $np_0 \ge 10$ and $n(1-p_0) \ge 10$.
    \item \textbf{Difference of two proportions} (Chapter 6.2): a two-outcome variable compared between two groups. Test whether the proportion is the same in both groups.\\
    \emph{Example:} Is the survival rate different for female and male Titanic passengers?\\
    \emph{Conditions:} independence within each group and between the groups; success-failure in each group using the pooled proportion $\hat{p}_{pool}$ = (total successes in both groups)$/(n_1 + n_2)$: at least 10 expected successes and 10 expected failures in each group.
    \item \textbf{Chi-square goodness-of-fit} (Chapter 6.3): one categorical variable with three or more outcomes. Test whether the counts match a hypothesized distribution.\\
    \emph{Example:} Are births equally likely on all seven days of the week?\\
    \emph{Conditions:} independence; every expected count is at least 5; degrees of freedom $> 1$ (at least 3 categories).
    \item \textbf{Chi-square test of independence} (Chapter 6.4): two categorical variables, at least one with three or more outcomes. Test whether the two variables are associated.\\
    \emph{Example:} Is survival associated with passenger class (1st, 2nd, 3rd) on the Titanic?\\
    \emph{Conditions:} independence; every expected count is at least 5, where expected count = (row total)(column total)/(table total).
\end{itemize}

\subsubsection*{Numerical Data Analysis (choose at least one)}
For the $t$-tests below, the \textbf{normality} condition depends on the sample size: if $n < 30$, the data should be nearly normal with no clear outliers; if $n \ge 30$, there should be no extreme skew. Check it with a histogram or boxplot.
\begin{itemize}
    \item \textbf{One-sample mean} (Chapter 7.1): one numerical variable. Compare its mean to a hypothesized value $\mu_0$.\\
    \emph{Example:} Do college students sleep less than 8 hours per night on average?\\
    \emph{Conditions:} independence; normality of the data.
    \item \textbf{Paired means} (Chapter 7.2): two numerical measurements on the same individuals (e.g., before and after). Test whether the mean of the differences is zero.\\
    \emph{Example:} Do students score differently on the final exam than on the midterm, on average?\\
    \emph{Conditions:} the pairs are independent of each other; normality of the \emph{differences}.
    \item \textbf{Difference of two means} (Chapter 7.3): one numerical variable compared between two independent (not paired) groups. Test whether the two group means are equal.\\
    \emph{Example:} Is the average birth weight different for babies of mothers who smoke and mothers who don't?\\
    \emph{Conditions:} independence within each group and between the groups; normality in \emph{each} group.
    \item \textbf{ANOVA} (Chapter 7.5): one numerical variable compared across three or more groups. Test whether all the group means are equal.\\
    \emph{Example:} Does average gas mileage differ for cars with 4, 6, and 8 cylinders?\\
    \emph{Conditions:} independence within and between groups; the observations in each group are nearly normal (especially important for small groups); constant variance: the groups have roughly equal spread (compare the standard deviations or side-by-side boxplots).
\end{itemize}

\noindent\textbf{What about regression?} Questions about how one numerical variable changes with another numerical variable are answered with \textbf{linear regression}. Regression is \textbf{not required} for this project (it is not taught until MA214). You may still show a scatterplot of two numerical variables as part of your EDA, but your hypothesis test should be one of the tests above.\\
\emph{Example:} On average, how much does a car's gas mileage change for each additional 1{,}000 pounds of weight? Is this value significantly different from zero? (This is a regression question, not a difference of means question.)


\end{document}
